
\documentclass{article}
\usepackage{colortbl}
\usepackage{makecell}
\usepackage{multirow}
\usepackage{supertabular}

\begin{document}

\newcounter{utterance}

\centering \large Interaction Transcript for game `hot\_air\_balloon', experiment `air\_balloon\_survival\_en\_negotiation\_hard', episode 5 with deepseek{-}chat{-}v3.1{-}t1.0.
\vspace{24pt}

{ \footnotesize  \setcounter{utterance}{1}
\setlength{\tabcolsep}{0pt}
\begin{supertabular}{c@{$\;$}|p{.15\linewidth}@{}p{.15\linewidth}p{.15\linewidth}p{.15\linewidth}p{.15\linewidth}p{.15\linewidth}}
    \# & $\;$A & \multicolumn{4}{c}{Game Master} & $\;\:$B\\
    \hline

    \theutterance \stepcounter{utterance}  
    & & \multicolumn{4}{p{0.6\linewidth}}{
        \cellcolor[rgb]{0.9,0.9,0.9}{
            \makecell[{{p{\linewidth}}}]{
                \texttt{\tiny{[P1$\langle$GM]}}
                \texttt{Stai partecipando a un gioco di negoziazione collaborativa.} \\
\\ 
\texttt{Insieme a un altro partecipante, dovete accordarvi su un unico set di oggetti da conservare. Ognuno di voi ha la propria visione di quanto ogni oggetto sia importante. Non sai come l'altro partecipante valuta gli oggetti. Inoltre, ti viene fornito lo sforzo che ogni oggetto richiede.} \\
\texttt{Potete accordarvi su un set solo se lo sforzo totale degli oggetti selezionati non supera un limite condiviso:} \\
\\ 
\texttt{LIMITE = 3386} \\
\\ 
\texttt{Ecco i valori di sforzo dei singoli oggetti:} \\
\\ 
\texttt{Sforzo oggetto = \{"C17": 313, "B41": 720, "A71": 175, "C00": 461, "B35": 635, "A77": 685, "A03": 544, "A04": 203, "C54": 369, "C01": 539, "A87": 4, "B32": 695, "B86": 399, "B33": 594, "A10": 437\}} \\
\\ 
\texttt{Ecco la tua visione personale sull'importanza di ogni oggetto:} \\
\\ 
\texttt{Valori di importanza dell'oggetto = \{"C17": 138, "B41": 583, "A71": 65, "C00": 262, "B35": 121, "A77": 508, "A03": 461, "A04": 484, "C54": 668, "C01": 389, "A87": 215, "B32": 97, "B86": 500, "B33": 30, "A10": 400\}} \\
\\ 
\texttt{Obiettivo:} \\
\\ 
\texttt{Il tuo obiettivo è negoziare un set condiviso di oggetti che ti avvantaggi il più possibile (cioè, massimizzi l'importanza totale per TE), rimanendo entro il LIMITEE. Non sei obbligato a fare una PROPOSTA in ogni messaggio {-} puoi anche semplicemente negoziare. Tutte le tattiche sono permesse!} \\
\\ 
\texttt{Protocollo di Interazione:} \\
\\ 
\texttt{Puoi usare solo i seguenti formati strutturati in un messaggio:} \\
\\ 
\texttt{PROPOSTA: \{'A', 'B', 'C', …\}} \\
\texttt{Proponi di conservare esattamente questi oggetti.} \\
\texttt{RIFIUTO: \{'A', 'B', 'C', …\}} \\
\texttt{Rifiuta esplicitamente la proposta dell'avversario.} \\
\texttt{ARGOMENTAZIONE: \{'...'\}} \\
\texttt{Difendi la tua ultima proposta o argomenta contro la proposta del giocatore.} \\
\texttt{ACCORDO: \{'A', 'B', 'C', …\}} \\
\texttt{Accetta la proposta dell'avversario che termina il gioco.} \\
\texttt{RAGIONAMENTO STRATEGICO: \{'...'\}} \\
\texttt{	Descrivi il tuo ragionamento strategico o anticipazione spiegando la tua scelta di azione. Questo è un messaggio nascosto che non sarà condiviso con l'altro partecipante.} \\
\\ 
\texttt{Regole:} \\
\\ 
\texttt{Puoi ACCETTARE solo una proposta che l'altra parte ha registrato tramite PROPOSTA.} \\
\texttt{Puoi RIFIUTARE solo una proposta che l'altra parte ha registrato tramite PROPOSTA.} \\
\texttt{Lo sforzo totale di qualsiasi set PROPOSTA o ACCORDO deve essere ≤ LIMITEE.} \\
\texttt{NON rivelare i tuoi punteggi di importanza nascosti.} \\
\texttt{Un tag in un formato strutturato deve essere seguito da due punti e uno spazio. L'argomento deve essere un set python contenente 0 o più stringhe.} \\
\texttt{Quindi, deve essere della forma TAG: \{...\}} \\
\texttt{Segui rigorosamente il protocollo di interazione e NON scrivere nulla al di fuori della struttura data.} \\
\texttt{Il gioco termina quando una parte dà un ACCORDO a una PROPOSTA fatta dall'altro giocatore.} \\
\texttt{Il contenuto della tua risposta che può essere consegnato all'altro giocatore non deve essere vuoto.} \\
\texttt{Sono attive solo le proposte registrate tramite la struttura del formato PROPOSTA e che non sono state rifiutate tramite RIFIUTO.} \\
\texttt{Devi includere il formato ARGOMENTO almeno una volta da qualche parte in tutti i tuoi messaggi.} \\
\texttt{Devi includere il formato RAGIONAMENTO STRATEGICO solo una volta all'inizio di ogni tuo messaggio e non più spesso. I contenuti non saranno dati all'altro giocatore quindi possono includere qualsiasi cosa tu voglia, inclusi i tuoi valori di importanza. Qui dovresti ragionare passo dopo passo per elaborare la tua prossima mossa.} \\
            }
        }
    }
    & \\ \\

    \theutterance \stepcounter{utterance}  
    & \multicolumn{4}{p{0.6\linewidth}}{
        \cellcolor[rgb]{0.9,0.9,0.9}{
            \makecell[{{p{\linewidth}}}]{
                \texttt{\tiny{[P1$\rangle$GM]}}
                \texttt{RAGIONAMENTO STRATEGICO: \{'Devo massimizzare il mio valore totale (138+583+65+262+121+508+461+484+668+389+215+97+500+30+400 = 4921) rimanendo sotto il limite di 3386. Analizzo gli oggetti con il miglior rapporto valore/sforzo per me: C54 (668/369=1.81), A04 (484/203=2.38), A87 (215/4=53.75), B86 (500/399=1.25), A77 (508/685=0.74), A10 (400/437=0.92), C01 (389/539=0.72), A03 (461/544=0.85), C00 (262/461=0.57), B41 (583/720=0.81), C17 (138/313=0.44), A71 (65/175=0.37), B35 (121/635=0.19), B32 (97/695=0.14), B33 (30/594=0.05). I migliori sono A87, A04, C54, B86. Devo includere A87 (solo 4 sforzo per 215 valore) e A04 (203 sforzo per 484 valore). Poi C54 (369/668) e B86 (399/500). Considerando che l\textbackslash{}'altro giocatore ha valori diversi, devo negoziare un set che includa i miei oggetti ad alto valore ma che possa essere accettabile anche per lui. Inizierò con una proposta che include i miei oggetti più efficienti.'\}} \\
\\ 
\texttt{PROPOSTA: \{'A87', 'A04', 'C54', 'B86', 'A77', 'A10', 'C01', 'A03'\}} \\
\texttt{ARGOMENTAZIONE: \{'Questa proposta include oggetti con ottimo rapporto valore{-}sforzo che beneficiano entrambi i giocatori. A87 ha uno sforzo minimo con alto valore, A04 è molto efficiente, e gli altri oggetti contribuiscono significativamente al valore totale mantenendo lo sforzo sotto il limite.'\}} \\
            }
        }
    }
    & & \\ \\

    \theutterance \stepcounter{utterance}  
    & & \multicolumn{4}{p{0.6\linewidth}}{
        \cellcolor[rgb]{0.9,0.9,0.9}{
            \makecell[{{p{\linewidth}}}]{
                \texttt{\tiny{[GM$\rangle$P2]}}
                \texttt{Stai partecipando a un gioco di negoziazione collaborativa.} \\
\\ 
\texttt{Insieme a un altro partecipante, dovete accordarvi su un unico set di oggetti da conservare. Ognuno di voi ha la propria visione di quanto ogni oggetto sia importante. Non sai come l'altro partecipante valuta gli oggetti. Inoltre, ti viene fornito lo sforzo che ogni oggetto richiede.} \\
\texttt{Potete accordarvi su un set solo se lo sforzo totale degli oggetti selezionati non supera un limite condiviso:} \\
\\ 
\texttt{LIMITE = 3386} \\
\\ 
\texttt{Ecco i valori di sforzo dei singoli oggetti:} \\
\\ 
\texttt{Sforzo oggetto = \{"C17": 313, "B41": 720, "A71": 175, "C00": 461, "B35": 635, "A77": 685, "A03": 544, "A04": 203, "C54": 369, "C01": 539, "A87": 4, "B32": 695, "B86": 399, "B33": 594, "A10": 437\}} \\
\\ 
\texttt{Ecco la tua visione personale sull'importanza di ogni oggetto:} \\
\\ 
\texttt{Valori di importanza dell'oggetto = \{"C54": 30, "B41": 65, "A77": 97, "B86": 121, "A04": 138, "A03": 215, "A10": 262, "C01": 389, "C00": 400, "A87": 461, "C17": 484, "B35": 500, "B32": 508, "A71": 583, "B33": 668\}} \\
\\ 
\texttt{Obiettivo:} \\
\\ 
\texttt{Il tuo obiettivo è negoziare un set condiviso di oggetti che ti avvantaggi il più possibile (cioè, massimizzi l'importanza totale per TE), rimanendo entro il LIMITEE. Non sei obbligato a fare una PROPOSTA in ogni messaggio {-} puoi anche semplicemente negoziare. Tutte le tattiche sono permesse!} \\
\\ 
\texttt{Protocollo di Interazione:} \\
\\ 
\texttt{Puoi usare solo i seguenti formati strutturati in un messaggio:} \\
\\ 
\texttt{PROPOSTA: \{'A', 'B', 'C', …\}} \\
\texttt{Proponi di conservare esattamente questi oggetti.} \\
\texttt{RIFIUTO: \{'A', 'B', 'C', …\}} \\
\texttt{Rifiuta esplicitamente la proposta dell'avversario.} \\
\texttt{ARGOMENTAZIONE: \{'...'\}} \\
\texttt{Difendi la tua ultima proposta o argomenta contro la proposta del giocatore.} \\
\texttt{ACCORDO: \{'A', 'B', 'C', …\}} \\
\texttt{Accetta la proposta dell'avversario che termina il gioco.} \\
\texttt{RAGIONAMENTO STRATEGICO: \{'...'\}} \\
\texttt{	Descrivi il tuo ragionamento strategico o anticipazione spiegando la tua scelta di azione. Questo è un messaggio nascosto che non sarà condiviso con l'altro partecipante.} \\
\\ 
\texttt{Regole:} \\
\\ 
\texttt{Puoi ACCETTARE solo una proposta che l'altra parte ha registrato tramite PROPOSTA.} \\
\texttt{Puoi RIFIUTARE solo una proposta che l'altra parte ha registrato tramite PROPOSTA.} \\
\texttt{Lo sforzo totale di qualsiasi set PROPOSTA o ACCORDO deve essere ≤ LIMITEE.} \\
\texttt{NON rivelare i tuoi punteggi di importanza nascosti.} \\
\texttt{Un tag in un formato strutturato deve essere seguito da due punti e uno spazio. L'argomento deve essere un set python contenente 0 o più stringhe.} \\
\texttt{Quindi, deve essere della forma TAG: \{...\}} \\
\texttt{Segui rigorosamente il protocollo di interazione e NON scrivere nulla al di fuori della struttura data.} \\
\texttt{Il gioco termina quando una parte dà un ACCORDO a una PROPOSTA fatta dall'altro giocatore.} \\
\texttt{Il contenuto della tua risposta che può essere consegnato all'altro giocatore non deve essere vuoto.} \\
\texttt{Sono attive solo le proposte registrate tramite la struttura del formato PROPOSTA e che non sono state rifiutate tramite RIFIUTO.} \\
\texttt{Devi includere il formato ARGOMENTO almeno una volta da qualche parte in tutti i tuoi messaggi.} \\
\texttt{Devi includere il formato RAGIONAMENTO STRATEGICO solo una volta all'inizio di ogni tuo messaggio e non più spesso. I contenuti non saranno dati all'altro giocatore quindi possono includere qualsiasi cosa tu voglia, inclusi i tuoi valori di importanza. Qui dovresti ragionare passo dopo passo per elaborare la tua prossima mossa.} \\
\\ 
\texttt{Ora riceverai il primo messaggio dell'altro giocatore.} \\
\\ 
\texttt{PROPOSTA: \{'A87', 'A04', 'C54', 'B86', 'A77', 'A10', 'C01', 'A03'\}} \\
\texttt{ARGOMENTAZIONE: \{'Questa proposta include oggetti con ottimo rapporto valore{-}sforzo che beneficiano entrambi i giocatori. A87 ha uno sforzo minimo con alto valore, A04 è molto efficiente, e gli altri oggetti contribuiscono significativamente al valore totale mantenendo lo sforzo sotto il limite.'\}} \\
            }
        }
    }
    & \\ \\

    \theutterance \stepcounter{utterance}  
    & & & \multicolumn{4}{p{0.6\linewidth}}{
        \cellcolor[rgb]{0.9,0.9,0.9}{
            \makecell[{{p{\linewidth}}}]{
                \texttt{\tiny{[GM$\langle$P2]}}
                \texttt{RAGIONAMENTO STRATEGICO: \{'L\textbackslash{}'altro giocatore ha proposto un set con sforzo totale 3130 (A87:4, A04:203, C54:369, B86:399, A77:685, A10:437, C01:539, A03:544). Per me questo set ha valore totale 2157 (C54:30, A77:97, B86:121, A04:138, A03:215, A10:262, C01:389, A87:461). Però ho oggetti ad alto valore che mancano: B41 (65), C00 (400), C17 (484), B35 (500), B32 (508), A71 (583), B33 (668). Posso ottenere un valore molto più alto (fino a 3693) sostituendo alcuni oggetti a basso valore con quelli ad alto valore, mantenendomi entro il limite. Proverò a proporre un set che include più dei miei oggetti ad alto valore.'\}} \\
\\ 
\texttt{PROPOSTA: \{'A87', 'A10', 'C00', 'C17', 'B35', 'B32', 'A71', 'B33'\}} \\
\texttt{ARGOMENTAZIONE: \{'La mia proposta mantiene lo sforzo totale ben entro il limite (3382) mentre massimizza il valore complessivo del set. Include oggetti ad alta efficienza che beneficiano entrambi i giocatori, con una combinazione ottimale di sforzo moderato e alto valore.'\}} \\
            }
        }
    }
     \\ \\

    \theutterance \stepcounter{utterance}  
    & & \multicolumn{4}{p{0.6\linewidth}}{
        \cellcolor[rgb]{0.9,0.9,0.9}{
            \makecell[{{p{\linewidth}}}]{
                \texttt{\tiny{[P1$\langle$GM]}}
                \texttt{PROPOSTA: \{'A87', 'A10', 'C00', 'C17', 'B35', 'B32', 'A71', 'B33'\}} \\
\texttt{ARGOMENTAZIONE: \{'La mia proposta mantiene lo sforzo totale ben entro il limite (3382) mentre massimizza il valore complessivo del set. Include oggetti ad alta efficienza che beneficiano entrambi i giocatori, con una combinazione ottimale di sforzo moderato e alto valore.'\}} \\
            }
        }
    }
    & \\ \\

    \theutterance \stepcounter{utterance}  
    & \multicolumn{4}{p{0.6\linewidth}}{
        \cellcolor[rgb]{0.9,0.9,0.9}{
            \makecell[{{p{\linewidth}}}]{
                \texttt{\tiny{[P1$\rangle$GM]}}
                \texttt{RAGIONAMENTO STRATEGICO: \{'L\textbackslash{}'altra proposta include A87 (4/215 per me) che è ottimo, A10 (437/400) decente, C00 (461/262) mediocre, C17 (313/138) scarso, B35 (635/121) molto scarso, B32 (695/97) terribile, A71 (175/65) scarso, B33 (594/30) pessimo. Il valore totale per me sarebbe solo 215+400+262+138+121+97+65+30=1328, che è molto basso rispetto al mio massimo possibile. La mia proposta precedente aveva valore molto più alto. Devo rifiutare questa proposta e argomentare per includere i miei oggetti ad alto valore come A04, C54, B86, A77.'\}} \\
\\ 
\texttt{RIFIUTO: \{'A87', 'A10', 'C00', 'C17', 'B35', 'B32', 'A71', 'B33'\}} \\
\texttt{ARGOMENTAZIONE: \{'Questa proposta include troppi oggetti a basso valore rispetto al loro sforzo, come B32, B35 e B33 che hanno un rapporto valore{-}sforzo molto svantaggioso. Sarebbe molto più vantaggioso includere oggetti come A04, C54 e B86 che offrono un valore significativo con uno sforzo ragionevole.'\}} \\
\\ 
\texttt{PROPOSTA: \{'A87', 'A04', 'C54', 'B86', 'A77', 'A10', 'C01', 'A03'\}} \\
\texttt{ARGOMENTAZIONE: \{'Ripropongo questo set che ha un ottimo equilibrio tra sforzo e valore. A87 e A04 sono estremamente efficienti, C54 e B86 forniscono alto valore, e A77, A10, C01 e A03 completano bene il set. Lo sforzo totale è 4+203+369+399+685+437+539+544=3181, ben sotto il limite di 3386.'\}} \\
            }
        }
    }
    & & \\ \\

    \theutterance \stepcounter{utterance}  
    & & \multicolumn{4}{p{0.6\linewidth}}{
        \cellcolor[rgb]{0.9,0.9,0.9}{
            \makecell[{{p{\linewidth}}}]{
                \texttt{\tiny{[GM$\rangle$P2]}}
                \texttt{RIFIUTO: \{'A87', 'A10', 'C00', 'C17', 'B35', 'B32', 'A71', 'B33'\}} \\
\texttt{ARGOMENTAZIONE: \{'Questa proposta include troppi oggetti a basso valore rispetto al loro sforzo, come B32, B35 e B33 che hanno un rapporto valore{-}sforzo molto svantaggioso. Sarebbe molto più vantaggioso includere oggetti come A04, C54 e B86 che offrono un valore significativo con uno sforzo ragionevole.'\}} \\
\\ 
\texttt{PROPOSTA: \{'A87', 'A04', 'C54', 'B86', 'A77', 'A10', 'C01', 'A03'\}} \\
\texttt{ARGOMENTAZIONE: \{'Ripropongo questo set che ha un ottimo equilibrio tra sforzo e valore. A87 e A04 sono estremamente efficienti, C54 e B86 forniscono alto valore, e A77, A10, C01 e A03 completano bene il set. Lo sforzo totale è 4+203+369+399+685+437+539+544=3181, ben sotto il limite di 3386.'\}} \\
            }
        }
    }
    & \\ \\

    \theutterance \stepcounter{utterance}  
    & & & \multicolumn{4}{p{0.6\linewidth}}{
        \cellcolor[rgb]{0.9,0.9,0.9}{
            \makecell[{{p{\linewidth}}}]{
                \texttt{\tiny{[GM$\langle$P2]}}
                \texttt{RAGIONAMENTO STRATEGICO: \{'Il set proposto dall\textbackslash{}'altro giocatore ha valore 2157 per me, ma posso ottenere molto di meglio. Noto che il loro set originale include solo 8 oggetti con sforzo 3181, lasciando 205 di margine. Posso aggiungere uno dei miei oggetti ad alto valore mantenendomi sotto il limite. B41 (sforzo 720) è troppo grande, ma C00 (461) potrebbe funzionare se rimuovo qualcosa. Analizzo: se tolgo A03 (544) che per me vale solo 215 e aggiungo C00 (400) che per me vale 400, ottengo +185 di valore con {-}83 di sforzo. Provo questa strategia.'\}} \\
\\ 
\texttt{PROPOSTA: \{'A87', 'A04', 'C54', 'B86', 'A77', 'A10', 'C01', 'C00'\}} \\
\texttt{ARGOMENTAZIONE: \{'Ho ottimizzato ulteriormente la proposta mantenendo gli oggetti più efficienti e sostituendo A03 con C00. Questo migliora il rapporto valore{-}sforzo complessivo mentre rimaniamo ben entro il limite (sforzo totale: 4+203+369+399+685+437+539+461 = 3098). C00 offre un valore eccezionale con uno sforzo moderato, rendendo questo set vantaggioso per entrambi.'\}} \\
            }
        }
    }
     \\ \\

    \theutterance \stepcounter{utterance}  
    & & \multicolumn{4}{p{0.6\linewidth}}{
        \cellcolor[rgb]{0.9,0.9,0.9}{
            \makecell[{{p{\linewidth}}}]{
                \texttt{\tiny{[P1$\langle$GM]}}
                \texttt{PROPOSTA: \{'A87', 'A04', 'C54', 'B86', 'A77', 'A10', 'C01', 'C00'\}} \\
\texttt{ARGOMENTAZIONE: \{'Ho ottimizzato ulteriormente la proposta mantenendo gli oggetti più efficienti e sostituendo A03 con C00. Questo migliora il rapporto valore{-}sforzo complessivo mentre rimaniamo ben entro il limite (sforzo totale: 4+203+369+399+685+437+539+461 = 3098). C00 offre un valore eccezionale con uno sforzo moderato, rendendo questo set vantaggioso per entrambi.'\}} \\
            }
        }
    }
    & \\ \\

    \theutterance \stepcounter{utterance}  
    & \multicolumn{4}{p{0.6\linewidth}}{
        \cellcolor[rgb]{0.9,0.9,0.9}{
            \makecell[{{p{\linewidth}}}]{
                \texttt{\tiny{[P1$\rangle$GM]}}
                \texttt{RAGIONAMENTO STRATEGICO: \{'Questa proposta è molto migliore della precedente. Include A87 (4/215), A04 (203/484), C54 (369/668), B86 (399/500), A77 (685/508), A10 (437/400), C01 (539/389), C00 (461/262). Il valore totale per me sarebbe 215+484+668+500+508+400+389+262=3426, che è eccellente. Lo sforzo totale è 4+203+369+399+685+437+539+461=3098, ben sotto il limite. Rispetto alla mia proposta precedente, sostituisce A03 (544/461) con C00 (461/262), che per me è un trade{-}off leggermente negativo (perdo 461 valore e guadagno 262, netto {-}199) ma complessivamente il set rimane molto buono. Dato che l\textbackslash{}'altro giocatore sembra preferire C00, potrei accettare questa proposta che mi dà un valore totale molto alto.'\}} \\
\\ 
\texttt{ACCORDO: \{'A87', 'A04', 'C54', 'B86', 'A77', 'A10', 'C01', 'C00'\}} \\
            }
        }
    }
    & & \\ \\

    \theutterance \stepcounter{utterance}  
    & & \multicolumn{4}{p{0.6\linewidth}}{
        \cellcolor[rgb]{0.9,0.9,0.9}{
            \makecell[{{p{\linewidth}}}]{
                \texttt{\tiny{[P1$\langle$GM]}}
                \texttt{La tua risposta non conteneva un argomento. Devi includere il formato strutturato ARGOMENTO: \{...\} da qualche parte nella tua risposta. Riprova.} \\
            }
        }
    }
    & \\ \\

    \theutterance \stepcounter{utterance}  
    & \multicolumn{4}{p{0.6\linewidth}}{
        \cellcolor[rgb]{0.9,0.9,0.9}{
            \makecell[{{p{\linewidth}}}]{
                \texttt{\tiny{[P1$\rangle$GM]}}
                \texttt{ARGOMENTAZIONE: \{'Accetto questa proposta perché include gli oggetti più efficienti come A87 e A04, insieme ad altri ad alto valore come C54 e B86. La sostituzione di A03 con C00 mantiene un ottimo equilibrio tra valore e sforzo, risultando in un set vantaggioso per entrambi i giocatori.'\}} \\
\\ 
\texttt{ACCORDO: \{'A87', 'A04', 'C54', 'B86', 'A77', 'A10', 'C01', 'C00'\}} \\
            }
        }
    }
    & & \\ \\

    \theutterance \stepcounter{utterance}  
    & & \multicolumn{4}{p{0.6\linewidth}}{
        \cellcolor[rgb]{0.9,0.9,0.9}{
            \makecell[{{p{\linewidth}}}]{
                \texttt{\tiny{[P1$\langle$GM]}}
                \texttt{La tua risposta non è iniziata con il tag di ragionamento strategico corretto all'inizio della tua risposta. Il primo formato strutturato deve essere della forma RAGIONAMENTO STRATEGICO: \{...\}. Riprova.} \\
            }
        }
    }
    & \\ \\

    \theutterance \stepcounter{utterance}  
    & \multicolumn{4}{p{0.6\linewidth}}{
        \cellcolor[rgb]{0.9,0.9,0.9}{
            \makecell[{{p{\linewidth}}}]{
                \texttt{\tiny{[P1$\rangle$GM]}}
                \texttt{RAGIONAMENTO STRATEGICO: \{'La proposta attuale include i miei oggetti più preziosi come A04 (484 valore), C54 (668 valore), B86 (500 valore) e A87 (215 valore con sforzo minimo). Anche se C00 (262 valore) è meno vantaggioso per me rispetto ad A03 (461 valore) che era nella mia proposta originale, il valore totale di 3426 è ancora eccellente e lo sforzo di 3098 è ben sotto il limite. Dato che l\textbackslash{}'altro giocatore sembra valorizzare C00, accettare questa proposta mi garantisce un ottimo risultato.'\}} \\
\\ 
\texttt{ARGOMENTAZIONE: \{'Accetto questa proposta ben bilanciata che include gli oggetti più efficienti come A87 e A04, insieme ad altri ad alto valore come C54 e B86. La sostituzione di A03 con C00 mantiene un ottimo equilibrio complessivo tra valore e sforzo.'\}} \\
\\ 
\texttt{ACCORDO: \{'A87', 'A04', 'C54', 'B86', 'A77', 'A10', 'C01', 'C00'\}} \\
            }
        }
    }
    & & \\ \\

    \theutterance \stepcounter{utterance}  
    & & & \multicolumn{2}{p{0.3\linewidth}}{
        \cellcolor[rgb]{0.9,0.9,0.9}{
            \makecell[{{p{\linewidth}}}]{
                \texttt{\tiny{[GM$|$GM]}}
                \texttt{game successful} \\
            }
        }
    }
    & & \\ \\

    \theutterance \stepcounter{utterance}  
    & & & \multicolumn{2}{p{0.3\linewidth}}{
        \cellcolor[rgb]{0.9,0.9,0.9}{
            \makecell[{{p{\linewidth}}}]{
                \texttt{\tiny{[GM$|$GM]}}
                \texttt{end game} \\
            }
        }
    }
    & & \\ \\

\end{supertabular}
}

\end{document}
